% Cómo leer este apunte
\section{Cómo leer este apunte}\label{cap:00}

El documento está organizado en capas que permiten volver sobre las mismas ideas con una
profundidad creciente:

\begin{center}
\begin{tabularx}{\textwidth}{>{\bfseries}p{34mm} X}
\toprule
Recorrido & Objetivo\\
\midrule
Primera lectura & Comprender qué representa cada concepto, qué se mueve y qué permanece.
Incluye una primera expresión algebraica, explicada término por término.\\
Segunda lectura & Construir figuras, seguir puntos e imágenes y reconocer visualmente
invariantes, centros, ejes y direcciones.\\
Tercera lectura & Relacionar las transformaciones con superficies, cartografía y problemas
de Agrimensura.\\
Apéndices & Reunir matrices, demostraciones, desarrollos de Taylor, casos excepcionales y
soluciones completas de las guías.\\
\bottomrule
\end{tabularx}
\end{center}

Dentro de cada capa, \textbf{dictado} identifica lo trabajado en clase,
\textbf{material docente} identifica los documentos enviados por la cátedra y
\textbf{desarrollo} reúne las explicaciones agregadas para hacer el material autónomo.

\emph{Consolidación 2026-09-21: hilo del profesor, rúbricas, bibliografía comentada.}
La sección~\ref{sec:hilo} ordena las clases según los resúmenes escritos por el profesor;
las cajas ``En palabras del profesor'' citan esos resúmenes; los apéndices~\ref{ap:rubricas},
\ref{ap:bibliografia} y~\ref{ap:notebooks} reúnen rúbricas de autoevaluación, la bibliografía
comentada y el material de clase.

\begin{idea}[Método de estudio propuesto por el docente]
Para cada definición o teorema conviene producir:
\begin{enumerate}
  \item una construcción o dibujo que muestre qué puntos se mueven y cuáles quedan fijos;
  \item una fórmula vectorial;
  \item un ejemplo numérico que permita comprobar la conclusión;
  \item una frase que identifique la propiedad esencial: distancia, ángulo, paralelismo,
  orientación o punto fijo.
\end{enumerate}
\end{idea}

\subsection{Desaprender el CBC: del álgebra de pizarrón al análisis geométrico del terreno}\label{sec:desaprender}

Encuadre propio, de alumno a alumno. El álgebra y el análisis del CBC dejan reflejos que
funcionan en el pizarrón y que esta materia pide cambiar, no porque estén mal sino porque
acá los objetos tienen un significado en el terreno. Son cuatro cambios de chip; la
Estación 0 del sitio de la materia (\texttt{00\_desaprender\_el\_cbc}, apéndice~\ref{ap:notebooks})
los trabaja con figuras, y la subsección~\ref{sec:puntos-formas-grillas} los resume en una
tabla.
\begin{enumerate}
  \item \textbf{Todo vector es un punto.} No una flecha abstracta: un mojón, o el
  desplazamiento entre dos mojones. Una matriz no es una tabla de números sino una manera
  de mover todos los mojones a la vez, o de cambiar la grilla con que se los mide.
  \item \textbf{El gradiente y las formas lineales no son vectores.} Son reglas graduadas
  tiradas sobre el plano (un mazo de rectas de nivel equiespaciadas) que miden vectores sin
  necesitar producto interno, y por eso se transforman del lado contrario, con \(A^{-1}\).
  El ``vector gradiente perpendicular a la curva de nivel'' existe recién cuando hay una
  métrica que lo levanta.
  \item \textbf{Las coordenadas son etiquetas; los metros los pone la métrica.} Latitud y
  longitud no miden distancias; la primera forma fundamental es la regla que convierte un
  paso en \((\varphi,\lambda)\) en metros sobre el terreno, y cambia de punto a punto.
  \item \textbf{Derivar es álgebra lineal en el parche tangente, y curvarse es una forma
  cuadrática.} La derivada de una parametrización es la jacobiana, una transformación
  lineal entre el plano de los parámetros y el plano tangente; lo que la superficie hace
  más allá de ese plano (cuánto se aparta en cada rumbo) lo dice una forma cuadrática, y
  diagonalizarla da las direcciones principales.
\end{enumerate}

\subsection{Convención para componer}

En este apunte, el producto \(F G\) significa ``primero aplicar \(G\) y después \(F\)'':
\[
  (F G)(x)=F\bigl(G(x)\bigr).
\]
Esta convención es importante porque, en general, cambiar el orden cambia el resultado.

\subsection{Mapa conceptual}

\begin{center}
\begin{tikzpicture}[
  node distance=7mm and 10mm,
  every node/.style={font=\small,align=center},
  box/.style={draw=teal,fill=teallight!40,rounded corners=2mm,
              minimum width=31mm,minimum height=9mm},
  arr/.style={-{Latex[length=2.4mm]},draw=tealdark,line width=.7pt}
]
  \node[box] (rp) {reflexión\\puntual};
  \node[box,right=of rp] (rr) {reflexión\\en una recta};
  \node[box,below=of rp] (tr) {traslación};
  \node[box,below=of rr] (rot) {rotación};
  \node[box,below=of tr] (hom) {homotecia};
  \node[box,below right=7mm and -5mm of hom] (sim) {similaridad};
  \draw[arr] (rp)--node[left]{dos centros}(tr);
  \draw[arr] (rr)--node[right]{ejes paralelos}(tr);
  \draw[arr] (rr)--node[right]{ejes secantes}(rot);
  \draw[arr] (tr)--(sim);
  \draw[arr] (rot)--(sim);
  \draw[arr] (hom)--(sim);
\end{tikzpicture}
\end{center}

